\chapter{Εισαγωγή}\label{ch:intro}

\section{Κίνητρο}\label{sec:intro-motivation}
Το Διαδίκτυο των Πραγμάτων (Internet of Things, IoT) είναι μία από τις πλέον διαδεδομένες εμπορικές τεχνολογίες σήμερα. Πρόκειται για ένα οικοσύστημα διασυνδεδεμένων «έξυπνων» συσκευών (λαμπτήρων, καμερών, θερμοστατών, αισθητήρων κ.ά.) που έχουν την δυνατότητα να αυτοματοποιούνται με βάση τις ανάγκες και τις συνήθειες κάθε χρήστη, επιτελώντας έναν προσωποποιημένο σκοπό. Η αυτοματοποίηση αυτή γίνεται μέσω κανόνων της μορφής «αν η κάμερα εντοπίσει κίνηση, τότε άναψε το φως», οι οποίοι μετατρέπουν τις συμβατικές συσκευές σε ένα συνεκτικό, έξυπνο περιβάλλον. Ωστόσο, η δημιουργία τέτοιων κανόνων παραμένει μια χειροκίνητη διαδικασία, την οποία πολλοί χρήστες δεν ολοκληρώνουν ποτέ, με αποτέλεσμα να μην αξιοποιείται το σημαντικό πλεονέκτημα των έξυπνων συσκευών έναντι των συμβατικών.

Το πρόβλημα αυτό αποτελεί αντικείμενο των συστημάτων σύστασης (recommender systems), μεθόδων που έχουν την δυνατότητα να προτείνουν προσωποποιημένους κανόνες σε κάθε χρήστη, αναλύοντας στατιστικά τις υπάρχουσες προτιμήσεις του ίδιου, καθώς και του συνόλου των χρηστών. Ωστόσο, η εκπαίδευση ενός τέτοιου συστήματος απαιτεί δεδομένα εξ ορισμού ευαίσθητα: η λίστα των συσκευών ενός σπιτιού και οι αυτοματισμοί του αποκαλύπτουν το ωράριο, τις συνήθειες και την τοποθεσία των χρηστών. Η συλλογή αυτών των δεδομένων και η εκπαίδευση ενός μοντέλου σε έναν κεντρικό server (η κλασική δηλαδή προσέγγιση της Μηχανικής Μάθησης) εγείρει σοβαρούς κινδύνους για την ιδιωτικότητα των χρηστών. Η Ομοσπονδιακή Μάθηση (Federated Learning) προσφέρει μια εναλλακτική: την εκπαίδευση ενός μοντέλου απευθείας στην συσκευή κάθε χρήστη, χωρίς τα δεδομένα καθαυτά να εγκαταλείπουν την συσκευή \cite{mcmahan_communication-efficient_2017}. Το σημείο σύνδεσης της σύστασης κανόνων IoT με την ομοσπονδιακή μάθηση αποτελεί το κεντρικό αντικείμενο της παρούσας εργασίας.

\section{Αντικείμενο και στόχοι}\label{sec:intro-goal}
Αφορμή για την εργασία αποτέλεσε ο ανοιχτός διαγωνισμός σύστασης κανόνων της εταιρείας Wyze \cite{noauthor_wyzelabsrulerecommendation_nodate}, ο οποίος δημοσιοποίησε ένα εκτενές, πραγματικό σύνολο δεδομένων από συσκευές και κανόνες εκατοντάδων χιλιάδων χρηστών. Η λύση της δεύτερης θέσης του διαγωνισμού \cite{conti_conti748wyze-rule-recommendation_2024}, που χρησιμοποιείται εδώ ως σημείο αναφοράς, αντιμετωπίζει το πρόβλημα ως μια centralized πρόβλεψη συνδέσεων σε γράφους. Παρότι πετυχαίνει ισχυρές επιδόσεις, αγνοεί εντελώς το πρόβλημα της ιδιωτικότητας που καθιστά το πρόβλημα ενδιαφέρον σε ρεαλιστικές συνθήκες.

Στόχος της παρούσας εργασίας είναι να επαναδιατυπώσει το πρόβλημα σε federated πλαίσιο και να απαντήσει σε δύο ερωτήματα: (α) μπορεί ένα μοντέλο εκπαιδευμένο ομοσπονδιακά να προσεγγίσει ή και να ξεπεράσει την centralized λύση αναφοράς; και (β) πώς μεταβάλλεται η επίδοση ανάλογα με τον τύπο της ενορχήστρωσης: από την centralized εκπαίδευση, στη συνεργασία λίγων οργανισμών (cross-silo), έως το ρεαλιστικό σενάριο εκατομμυρίων μεμονωμένων συσκευών (cross-device);

Οι κύριες συνεισφορές της εργασίας είναι οι εξής:
\begin{itemize}
  \item Επαναδιατύπωση της σύστασης κανόνων IoT ως ομοσπονδιακό task, με μια ενιαία υλοποίηση που εκτελείται αυτούσια σε όλους τους τύπους ενορχήστρωσης.
  \item Τρεις εκδόσεις μοντέλου αυξανόμενης πολυπλοκότητας, με την τελευταία (CompGCN + ComplEx + InfoNCE) να βελτιώνει την centralized λύση αναφοράς κατά 71\% στη μετρική MRR.
  \item Μελέτη της επίδοσης του συστήματος σε τέσσερα σενάρια ενορχήστρωσης, η οποία καταδεικνύει ότι η βέλτιστη επιλογή μοντέλου εξαρτάται από τον τύπο της ενορχήστρωσης: το ισχυρότερο μοντέλο κυριαρχεί παντού εκτός από το πιο κατανεμημένο σενάριο, όπου ένα απλούστερο μοντέλο αποδεικνύεται καλύτερο.
  \item Ένα αυστηρό πρωτόκολλο αξιολόγησης cold-start με πλήρη leave-one-out μηχανισμό, που μετράει την πραγματική ικανότητα γενίκευσης σε άγνωστους χρήστες, πέρα από τη μετρική του διαγωνισμού.
\end{itemize}

\section{Related work}\label{sec:intro-related}
Η παρούσα εργασία στηρίζεται στην δημοσίευση της Wyze που συνόδευε τον διαγωνισμό \cite{yao_fedrule_2022,yao_yh-yaofedrule_2025}, η οποία εισάγει ένα σύστημα σύστασης κανόνων IoT βασισμένο σε federated GNNs. Το μοντέλο της Wyze μοντελοποιεί κάθε χρήστη ως έναν γράφο συσκευών και κανόνων και εκπαιδεύει ομοσπονδιακά ένα GNN για πρόβλεψη συνδέσεων. Η παρούσα διπλωματική μοιράζεται αυτή την δομή, αλλά την εμπλουτίζει, μελετώντας διαφοροποιήσεις στην απόδοση ανάλογα με το περιβάλλον federation, και εισάγει μια νέα relational αρχιτεκτονική για την βελτιστοποίηση του συστήματος.

Ευρύτερα, η εργασία απασχολεί τρία ερευνητικά πεδία: τα συστήματα σύστασης, την ομοσπονδιακή μάθηση και την βαθιά μάθηση σε γράφους, καθένα από τα οποία αναπτύσσεται στο \autoref{ch:theory}. Το ίδιο το πρόβλημα, ο διαγωνισμός της Wyze και η λύση αναφοράς περιγράφονται αναλυτικά στο \autoref{ch:problem}.

\section{Δομή της εργασίας}\label{sec:intro-structure}
Η εργασία οργανώνεται ως εξής: Το \autoref{ch:theory} παρουσιάζει το θεωρητικό υπόβαθρο: Μηχανική Μάθηση, Ομοσπονδιακή Μάθηση, θεωρία γράφων, GNNs και συστήματα σύστασης. Το \autoref{ch:problem} περιγράφει το πρόβλημα, τον διαγωνισμό της Wyze, τη μοντελοποίηση των δεδομένων ως γράφων και την λύση της δεύτερης θέσης. Το \autoref{ch:methodology} αναπτύσσει τη μεθοδολογία και τις τρεις εκδόσεις μοντέλου. Το \autoref{ch:experiments} παρουσιάζει την πειραματική διάταξη, την ανάλυση των δεδομένων, τα σενάρια ενορχήστρωσης, το πρωτόκολλο αξιολόγησης και τα αποτελέσματα. Τέλος, το \autoref{ch:conclusions} συνοψίζει τα συμπεράσματα, συζητά τους περιορισμούς και προτείνει μελλοντικές επεκτάσεις.
